\documentclass[11pt]{article}
\usepackage[margin=1in]{geometry}
\usepackage{booktabs}
\usepackage{tabularx}
\usepackage{float}
\usepackage{array}
\usepackage{xurl}
\usepackage{hyperref}
\emergencystretch=2em
\usepackage{xcolor}
\usepackage{enumitem}
\newcommand{\todo}[1]{\textcolor{red}{[TODO: #1]}}

\title{DS-GA 1011 Assignment 4: Fine-tuning Language Models}
\author{Rishit Tribhovanbhai Maheshwari (rm7336)}
\date{Fall 2026}

\begin{document}
\maketitle

\noindent\textbf{Code:} \url{https://github.com/Reemostat/nyu-dsga1011-a4}\\
\textbf{Model checkpoint (Q7):} \url{https://drive.google.com/drive/folders/1r81kBF-5GTwNHe0AyCS-G5l7I_yAfrBk}\\
(\texttt{t5\_best\_checkpoint.tar.gz}: fine-tuned \texttt{t5-small}, best epoch)

%==================================================================
\section*{Part 1: Fine-tuning BERT for Sentiment Classification}

\subsection*{Q1. Fine-tuning BERT}
I implemented the training loop in \texttt{do\_train}. For each of the 3 epochs and each batch, the loop
moves the batch to the GPU and runs a forward pass. The batch includes \texttt{labels}, so the model returns the cross-entropy loss.
The loop then calls \texttt{optimizer.zero\_grad()}, back-propagates, and steps both the AdamW optimizer (lr $5\times10^{-5}$) and the linear
learning-rate scheduler. All hyperparameters are the defaults: batch size 8, 3 epochs, \texttt{bert-base-cased}.

\textbf{Result:} fine-tuned on the full IMDB training set (25k reviews), the model reaches
\textbf{93.12\%} accuracy on the IMDB test set (25k reviews). Predictions are in \texttt{out\_original.txt}.

\subsection*{Q2a. Transformation design}
My transformation simulates a \emph{careless user who paraphrases loosely and types quickly}. It combines
synonym substitution with keyboard typos.

\begin{enumerate}[leftmargin=*]
  \item \textbf{Segmentation.} The review is split with the regex \verb/([^A-Za-z]+)/, keeping the delimiters. This gives an
  alternating sequence of alphabetic words and non-letter runs (spaces, punctuation, digits, HTML such as
  \texttt{<br />}). Only the alphabetic pieces can change, and concatenating the pieces rebuilds the text exactly.
  \item \textbf{Eligibility.} A word is eligible if it has at least 4 letters and is not in a small protected list.
  The list holds negations (\emph{not, never, nothing, without, cannot, \ldots}), common function words
  (\emph{this, that, with, from, have, very, \ldots}), contraction stems (\emph{didn, wasn, don, \ldots}) and the HTML token
  \emph{br}. The protected list keeps the label unchanged: WordNet would otherwise map \emph{not} to \emph{non}, for example.
  \item \textbf{Synonym replacement ($p=0.30$).} For each eligible word, draw $r\sim U(0,1)$. If $r<0.30$, collect
  every lemma of every WordNet synset of the lower-cased word, drop the word itself and non-ASCII lemmas, and replace
  underscores with spaces. Pick one candidate uniformly at random and copy the original capitalization
  (all-caps or first-letter capital). If WordNet has no synonym, the word stays unchanged.
  \item \textbf{Keyboard typo ($p=0.20$).} Otherwise, if $0.30\le r<0.50$, pick one letter of the word uniformly at
  random and replace it with a uniformly chosen neighbouring key on a QWERTY keyboard
  (e.g.\ \texttt{a}$\to$\{\texttt{q,w,s,z}\}), keeping its case. Each word gets at most one typo.
\end{enumerate}

\noindent\textbf{Example.} \emph{``I really loved this movie. The acting was superb and the story kept me engaged
until the very end.''} $\to$ \emph{``I really loved this mobie. The acging was siperb and the taradiddle keep me
engaged until the very end.''}

\noindent\textbf{Why it is reasonable.} Real users write reviews fast, often on phones, so slipped keys
(\emph{mobie}, \emph{siperb}) are common in user-generated text. Different users also describe the same opinion
with different words: \emph{film} vs.\ \emph{movie}, \emph{superb} vs.\ \emph{splendid}. A human reading the
transformed review still assigns the same sentiment. Negations and function words are protected, every edit is local
(one word), and the overall structure and sentiment-bearing phrases stay recognisable.

\subsection*{Q2b. Results on the transformed test set}
\begin{center}
\begin{tabular}{lc}
\toprule
Test set & Accuracy \\
\midrule
Original & 93.12\% \\
Transformed & 87.17\% \\
\midrule
Drop & \textbf{5.96 points} \\
\bottomrule
\end{tabular}
\end{center}
Predictions are in \texttt{out\_transformed.txt}.

\subsection*{Q3. Data augmentation}
I sampled 5{,}000 random training reviews (\texttt{shuffle(seed=42)}) and applied the same transformation
(\texttt{custom\_transform})
to them. I concatenated them with the full original training set (30k examples total) and trained a fresh
\texttt{bert-base-cased} model with the same hyperparameters.

\begin{center}
\begin{tabular}{lcc}
\toprule
Model & Original test & Transformed test \\
\midrule
Trained on original data (Q1/Q2) & 93.12\% & 87.17\% \\
Trained on augmented data (Q3)   & 93.03\% & 89.59\% \\
\midrule
Change & $-0.09$ & $+2.42$ \\
\bottomrule
\end{tabular}
\end{center}
Predictions are in \texttt{out\_augmented\_original.txt} and \texttt{out\_augmented\_transformed.txt}.

\textbf{(1) Transformed test data.} Performance improved by 2.42 points (87.17\%$\to$89.59\%), recovering about
40\% of the accuracy lost under the transformation.

\textbf{(2) Original test data.} Accuracy is essentially unchanged (93.12\%$\to$93.03\%, $-0.09$). A difference this
small is within run-to-run noise, so augmentation neither clearly helped nor hurt clean performance.

\textbf{Intuitive explanation.} During training, the original model only saw clean text. A typo such as
\emph{siperb} is broken by BERT's WordPiece tokenizer into rare sub-word pieces that the model never learned to tie
to sentiment, and an unusual synonym such as \emph{taradiddle} is effectively an unseen word. With augmentation,
the model sees noisy versions of reviews with the same label. It learns to rely on the remaining clean context and
to map corrupted sub-word sequences to the same meaning, which makes it more robust to this specific noise.
Clean reviews are still the majority of the training data (25k of 30k), so performance on clean text is preserved.
The gain is only partial because 5k noisy examples (one sixth of the data, a single noisy version per review)
cover only a small share of the possible typos and synonyms.

\textbf{Limitation.} Augmentation only helps for the shift you anticipated and simulated. Here the model became more
robust to \emph{my} typo/synonym noise, but that says nothing about other out-of-distribution shifts at test
time: different domains (product or restaurant reviews), sarcasm, other languages or dialects, paraphrase styles
WordNet cannot produce, or other typo patterns. In practice the true test-time distribution is unknown, so you
cannot enumerate every transformation to augment with. Synthetic noise can also differ in its statistics from
real noise, and aggressive augmentation can introduce label noise.

%==================================================================
\section*{Part 2: Fine-tuning T5 for Text-to-SQL}

\subsection*{Q4. Data statistics}
All lengths are counted in \texttt{google-t5/t5-small} tokenizer tokens. Vocabulary size is the number of distinct
token ids that occur in that split.

\begin{table}[H]
\centering
\begin{tabular}{lrr}
\toprule
Statistics Name & Train & Dev \\
\midrule
Number of examples & 4225 & 466 \\
Mean sentence length & 17.10 & 17.07 \\
Max sentence length & 59 & 43 \\
Mean SQL query length & 216.37 & 210.05 \\
Max SQL query length & 510 & 502 \\
Vocabulary size (natural language) & 791 & 465 \\
Vocabulary size (SQL) & 555 & 395 \\
\bottomrule
\end{tabular}
\caption{Data statistics before any pre-processing (T5 tokenizer, no special tokens).}
\end{table}

\begin{table}[H]
\centering
\begin{tabular}{lrr}
\toprule
Statistics Name & Train & Dev \\
\midrule
\multicolumn{3}{l}{\textbf{T5 fine-tuned (t5-small)}} \\
Mean sentence length & 23.10 & 23.07 \\
Max sentence length & 65 & 49 \\
Mean SQL query length & 217.57 & 211.24 \\
Max SQL query length & 511 & 503 \\
Vocabulary size (natural language) & 795 & 469 \\
Vocabulary size (SQL) & 556 & 396 \\
\bottomrule
\end{tabular}
\caption{Data statistics after pre-processing: task prefix on the input, \texttt{<}/\texttt{<=} placeholders in the
SQL, and the \texttt{</s>} EOS token appended to both sides.}
\end{table}

\subsection*{Q5. T5 fine-tuning details}
\begin{table}[H]
\centering
\begin{tabularx}{\textwidth}{lX}
\toprule
Design choice & Description \\
\midrule
Data processing &
The input becomes \texttt{"translate English to SQL: "} + the natural-language query. In the SQL, the comparison
operators \texttt{<=} and \texttt{<} are replaced by the words \texttt{LTE\_OP} and \texttt{LT\_OP} before tokenization,
because the T5 SentencePiece vocabulary has no \texttt{<} character (it would become \texttt{<unk>}). They are mapped
back after decoding, and runs of whitespace are collapsed. I verified that encoding then decoding reproduces all
4225 train and 466 dev SQL queries character-for-character. \\
Tokenization &
The default \texttt{T5TokenizerFast} (\texttt{google-t5/t5-small}) is used for both encoder and decoder, with
\texttt{</s>} appended to both sequences. No truncation is needed: the longest target is 511 tokens. The decoder input is the
target shifted right, with T5's decoder start token (the pad token, id 0) prepended. Batches are dynamically padded
with id 0, and the loss ignores pad positions. Decoding uses
\texttt{clean\_up\_tokenization\_spaces=False} so that the SQL spacing (e.g.\ \texttt{" , "}) is preserved. \\
Architecture &
The full pretrained \texttt{T5ForConditionalGeneration} (\texttt{t5-small}, 60.5M parameters) is fine-tuned:
encoder, decoder, shared embeddings and LM head. No layers are frozen. \\
Hyperparameters &
AdamW (lr $5\times10^{-4}$, $\beta=(0.9,0.999)$, $\epsilon=10^{-8}$, weight decay 0.01 except on biases and
LayerNorm), batch size 16, cosine schedule with 1 warm-up epoch and 20 scheduled epochs, token-level cross-entropy.
After every epoch the model is evaluated on dev by generating with beam search (4 beams, up to 512 new tokens) and
executing the SQL. The checkpoint with the best dev Record F1 is kept, with early stopping after 5 epochs without
improvement. The run hit the 3-hour job limit during epoch 15; the best checkpoint is from epoch 14 (index 13).
For the final predictions the saved checkpoint is decoded \emph{greedily} (up to 512 new tokens), which was
better than beam search on dev (Table~4). \\
\bottomrule
\end{tabularx}
\caption{Details of the best-performing T5 model configuration (fine-tuned).}
\end{table}

\subsection*{Q6. Results and error analysis}
\begin{table}[H]
\centering
\begin{tabular}{lcc}
\toprule
System & Query EM & F1 score \\
\midrule
\multicolumn{3}{l}{\textbf{Dev Results}} \\
\multicolumn{3}{l}{\emph{T5 fine-tuned}} \\
Full model, greedy decoding (final) & 67.60 & 82.91 \\
Full model, beam search (4 beams) & 63.09 & 80.16 \\
\midrule
\multicolumn{3}{l}{\textbf{Test Results}} \\
T5 fine-tuning (greedy) & 77.31$^*$ & 79.92 \\
\bottomrule
\end{tabular}
\caption{Development and test results (\%). Dev Query EM is exact match on the SQL string. $^*$The test leaderboard
cannot compute SQL EM (gold test SQL is hidden), so the test entry is Record EM.}
\end{table}

\noindent Dev Record F1 during training (beam 4):
epoch 1: 27.6, 2: 42.7, 3: 55.7, 4: 61.1, 5: 69.2, 6: 69.4, 7: 74.9, 8: 75.7, 9: 74.6, 10: 74.7, 11: 77.4,
12: 78.4, 13: 78.5, 14: 79.6.
Re-evaluating the saved best checkpoint gives Record F1 80.16, Record EM 77.25, SQL EM 63.09, with 7.94\% of
generated queries raising an SQL error (the small F1 difference from 79.6 comes from non-deterministic GPU batched beam search).

\noindent\textbf{Decoding ablation.} Decoding the same checkpoint greedily gives Record F1 82.91, Record EM 80.90,
SQL EM 67.60 and 7.08\% SQL errors, better than beam search on every metric. The dev loss is identical (0.0202),
so the only difference is the search procedure. A likely explanation: beam search ranks complete hypotheses by total log-probability, so it
favours shorter queries: hypotheses that drop a join or a \texttt{WHERE} clause pay fewer per-token penalties.
In this task, dropping one condition makes the records wrong. Greedy decoding stays on the token-by-token
most likely path, which the model was trained to follow with teacher forcing, and so produces the full,
long ATIS queries more reliably. The final system therefore uses greedy decoding.

\begin{table}[H]
\centering
\footnotesize
\setlength{\tabcolsep}{4pt}
\begin{tabularx}{\textwidth}{>{\raggedright\arraybackslash}p{1.9cm}>{\raggedright\arraybackslash}X>{\raggedright\arraybackslash}Xr}
\toprule
Error Type & Example of Error & Error Description & Statistics \\
\midrule
Wrong constant / value &
NL: \emph{``flights arriving in baltimore from philadelphia at about 4 o'clock''}.
Gold: \texttt{arrival\_time >= 1530 AND arrival\_time <= 1630}.
Predicted: \texttt{arrival\_time >= 1630 AND arrival\_time <= 1630}, which keeps only flights arriving exactly at 16:30. &
The SQL structure is right but a literal is wrong: a time window, date, number or code. Fuzzy expressions
(\emph{about 4}, \emph{morning}) must be mapped to numeric ranges, and the decoder has to copy or compute numbers
that appear in the NL only as words (\emph{four}, \emph{ninth}). &
48/466 \\
\midrule
Malformed SQL (syntax) &
Execution fails with \texttt{near ")": syntax error} (10), \texttt{incomplete input} (3) or an unterminated string
literal (1). &
The query has unbalanced parentheses or is cut off. Gold queries are long (mean $\approx$210 tokens, up to 510)
with nested \texttt{AND(}\ldots\texttt{)} groups and sub-queries. A sequence model keeps no explicit stack, so it
loses track of how many brackets are open. &
14/466 \\
\midrule
Undeclared table alias (schema) &
\texttt{no such column: flight\_fare\_1.flight\_id}: the WHERE clause joins on \texttt{flight\_fare\_1}, but
\texttt{flight\_fare flight\_fare\_1} never appears in the FROM list. Similar errors occur for \texttt{flight\_stop\_1},
\texttt{city\_2} and \texttt{airport\_1}. &
The FROM list and the WHERE conditions are produced far apart in the sequence, and nothing forces them to agree.
The model reproduces a familiar join chain without having declared every table it uses (8 cases, plus 1
ambiguous column name). &
9/466 \\
\midrule
Missing join condition (time-out) &
The query declares several tables in FROM without connecting all of them in WHERE, and execution hits the 120\,s
time-out. &
Tables that are declared but not connected make SQLite build a huge Cartesian product, so no records are returned.
This happens in the longest queries that mix flight and fare information. &
4/466 \\
\midrule
Wrong flight direction &
NL: \emph{``show me the united flights from bwi to denver''}. Gold: \texttt{from\_airport = \ldots 'BWI'} and
\texttt{to\_airport = \ldots 'DENVER'}. Predicted: \texttt{from\_airport} for \emph{both} BWI and Denver.
Also \emph{``i need to reverse the flight from pittsburgh to denver''}: the prediction keeps Pittsburgh as origin. &
All the right tables and constants are present, but a city is attached to the wrong column. Origin and destination
use identical join chains that differ only in \texttt{from\_airport}/\texttt{to\_airport}, so one wrong token
silently reverses or breaks the route. Cues like \emph{reverse} are rare in training. &
4/466 \\
\midrule
Wrong selected columns &
NL: \emph{``show me times for flights from san francisco to atlanta''}. Gold selects
\texttt{departure\_time, arrival\_time}; predicted selects only \texttt{departure\_time}. Also \emph{``define
airline us''}: only \texttt{airline\_code} instead of code, name and note. &
The WHERE clause is correct but the SELECT list is not. Which columns to return depends on vague words
(\emph{times}, \emph{define}, \emph{about}), and the training data is inconsistent about them. &
4/466 \\
\midrule
Wrong query intent / tables &
NL: \emph{``what is sa''}. Gold: day name for day code SA. Predicted: looks up airline code SA. Also
\emph{``flights into atlanta's airport''}: predicted \texttt{airport\_code = 'ATLANTA'} instead of joining through
the city table. &
The model queries the wrong part of the schema. These utterances are short or ambiguous, and ATIS was collected
as dialogues, so the context that would disambiguate them is missing. &
3/466 \\
\bottomrule
\end{tabularx}
\caption{Qualitative error analysis on the dev set for the final model (best checkpoint, greedy decoding).
Of 466 dev queries, 89 return records that differ from the gold records. 33 predictions raise an SQL error; 27 of
these are wrong, and the other 6 still count as correct because the gold query also returns no records.
Categories were assigned with \texttt{error\_analysis.py} (execution errors first, then literal mismatches, then
table-set mismatches) and the structural cases were checked by hand. The remaining 3 wrong-record cases are a
\texttt{COUNT} added to ``how many seats in a 734'', a sub-query missing its joins, and one long multi-stop query.}
\end{table}

\subsection*{Q7. Performance evaluation}
Test predictions (greedy decoding of the best checkpoint): \texttt{results/t5\_ft\_experiment\_test.sql} and \texttt{records/t5\_ft\_experiment\_test.pkl}.
Leaderboard results: Record F1 \textbf{79.92}, Record EM 77.31, with 43/432 queries raising an SQL error. This is above
the 65 F1 threshold. The test F1 is about 3 points below dev F1 (82.91), a small gap that is consistent with selecting the
checkpoint and the decoding method on dev.

\end{document}
